\documentclass[11pt]{article}
\usepackage[utf8]{inputenc}

%----- Configuración del estilo del documento------%
\usepackage{epsfig,graphicx}
\usepackage[left=2cm,right=2cm,top=1.8cm,bottom=2.3cm]{geometry}
\usepackage{fancyhdr}
\usepackage{lastpage}
\usepackage{url}
\usepackage{longtable}
\pagestyle{fancy}
\fancyhf{}
\fancyhead[R]{\ifnum\value{page}>1 \textbf{Teorema de Desargues}\fi}
\rfoot{\textit{Página \thepage}}

%Encabezado izquierdo: logo
\lfoot{\includegraphics[height=1.0cm]{Logo_FC.png}}



%------ Paquetes matemáticos básicos --------%
\usepackage{amsmath}
\usepackage{amssymb}
\usepackage{amsthm}
\usepackage[spanish]{babel}
\usepackage{graphicx}
\usepackage{hyperref}

\usepackage{tabularx}
\usepackage{xcolor}
\usepackage[table]{xcolor}
\usepackage{colortbl}
\usepackage{array, multirow, multicol, tabularx}
\usepackage{tcolorbox}
\newtheorem{theorem}{Theorem}[section]
\newtheorem{corollary}{Corollary}[theorem]
\newtheorem{lemma}[theorem]{Lemma}

%------si-------%
\definecolor{B}{HTML}{FFFFFF}
\definecolor{G}{HTML}{5e5e5e}
\definecolor{R2}{HTML}{d53d40}
\definecolor{A2}{HTML}{034190}
\definecolor{V2}{HTML}{7faa50}

\newcommand{\PrdI}[2]{\ensuremath{\langle #1, #2 \rangle}}
\newcommand{\R}{\mathbb{R}}
\newcommand{\C}{\mathbb{C}}
\newcommand{\N}{\mathbb{N}}
\newcommand{\Z}{\mathbb{Z}}
\newcommand{\Q}{\mathbb{Q}}
\renewcommand{\theenumi}{\Roman{enumi}}
\renewcommand{\labelenumi}{{\theenumi}.}

\begin{document}
	\makeatletter
        \renewenvironment{proof}[1][\proofname]{\par
            \pushQED{\qed}%
            \normalfont \topsep6\p@\@plus6\p@\relax
            \trivlist
            \item\relax
            {\itshape
            #1\@addpunct{.}}\par\vspace{\baselineskip}\ignorespaces
            }{%
            \popQED\endtrivlist\@endpefalse
            }
        \newenvironment{Solucion}[1][Soluci\'on]{\par
            \pushQED{\qed}%
            \normalfont \topsep6\p@\@plus6\p@\relax
            \trivlist
            \item\relax
            {\itshape
            #1\@addpunct{.}}\par\vspace{\baselineskip}\ignorespaces
            }{%
            \popQED\endtrivlist\@endpefalse
            }
            
    \makeatother

%------ Encabezado -------- %

\begin{center}
    \begin{minipage}{3cm}
    	\begin{center}
    		\includegraphics[height=3.4cm]{logo_unam.png}
    	\end{center}
    \end{minipage}\hfill
    \begin{minipage}{10cm}
    	\begin{center}
    	\textbf{\large Universidad Nacional Autónoma de México}\\[0.1cm]
        \textbf{Facultad de Ciencias}\\[0.1cm]
        \textbf{Ejercicios }\\[0.1cm]
         Teorema de Desargues\\[0.1cm]
         El\'ias L\'opez Rivera\\[0.1cm]
        \texttt{elias.lopezr,@ciencias.unam.mx }\\[0.1cm]
        Fecha:\,\,02/01/2026
    	\end{center}
    \end{minipage}\hfill
    \begin{minipage}{3cm}
    	\begin{center}
    		\includegraphics[height=3.4cm]{Logo_FC.png}
    	\end{center}
    \end{minipage}
\end{center}
\rule{17cm}{0.1mm}
\begin{tcolorbox}[title = \textcolor{black}{\textcolor{white}{Teorema de Desargues}}]
\textit{En un plano dos tri\'angulos $ABC$ y $A'B'C'$ estan en perspectiva desde un punto si y solo si estan en perspectiva desde una recta.\,\\
\begin{center}
    \includegraphics[width=0.5\textwidth]{Desargues.png}
\end{center}
}
\end{tcolorbox}

\begin{tcolorbox}[
	title = \textcolor{black}{\textcolor{white}{Problema 1}},]
\textit{¿Qu\'e eje es el de pespectiva de dos tri\'angulos homot\'eticos?}
\end{tcolorbox}

\begin{tcolorbox}[
	title = \textcolor{black}{\textcolor{white}{Problema 2}},]
\textit{Demostrar que si tres tri\'angulos est\'an en perspectiva respecto a un mismo punto,los tres ejes de perspectiva son cocurrentes.}
\end{tcolorbox}
\begin{proof}
    Tomemos los tri\'angulos $\Delta GEF$, $\Delta HIJ$ y $\Delta LMK$ con centro de perspectiva $A$, tomemos el eje de perspectiva $l_1$ de los tri\'angulos $\Delta GEF$ y $\Delta HIJ$, $l_2$ el eje de perspectiva de los tri\'angulos $\Delta HIJ$ y $\Delta LMK$, $l_3$ eje de perspectiva de los tri\'angulos $\Delta GEF$ y $\Delta LMK$,
    Tenemos que las rectas $\overline{EF}$ y $\overline{IJ}$ concurren con $l_1$ en $N$, ademas $\overline{EF}$ concurre con $\overline{LK}$ y $l_3$ en $Q$, ademas $\overline{LK}$ concurre con $l_2$ y $\overline{IH}$ en $T$, tenemos quue si se diera el caso de que cualesquiera dos de los puntos $N$, $Q$ y $T$ fueran iguales tendriamos necesariamente que los tres deben ser iguales, pues esto implicaria que las rectas $\overline{GF}$, $\overline{IH}$ y $\overline{LK}$ serian concurrentes en un solo punto, esto implicaria que $l_1$, $l_2$ y $l_3$ tamb\'ien serian concurrentes en el mismo punto, por tanto supongamos que los tres puntos son diferentes,
    ahora como consideramos que las tres rectas $\overline{GE}$, $\overline{HJ}$ y $\overline{LM}$ son diferentes, tenemos que los tres puntos no pueden ser colineales, pues esto implicaria que las rectas serian las mismas, esto quiere decir que $\Delta TQN$ es un tri\'angulo, si procedemos por un razonamiento an\'alogo sobre los puntos:
    \begin{itemize}
        \item $O=\overline{GF}\cap \overline{IJ}\cap l_1$
        \item $V=\overline{IJ}\cap \overline{ML}\cap l_2$
        \item $S=\overline{ML}\cap \overline{GF}\cap l_3$
    \end{itemize}
    Obtendremos el tri\'angulo $\Delta OVS$, afirmamos que los tri\'angulos $\Delta QVS$ y $\Delta TQN$, estan en perspectiva desde la recta $\overline{AF}$, pues sabemos que los puntos $I$, $L$ estan sobre esta recta y tenemosm las siguientes relaciones:
    \begin{itemize}
        \item los puntos $T$ , $N$ e $I$, $H$ son colineales, de la misma manera $O$, $I$ e $J$, $N$ lo son, por tanto $\overline{TN}$ y $\overline{OV}$ concurren en $I$, pues es el vertice del tri\'angulo $\Delta HIJ $
        \item Los puntos $E$, $F$ y $Q$, $N$ son colineales, de la misma manera $F$, $G$ y $O$, $S$ lo son, por tanto $\overline{QN}$ y $\overline{SO}$ concurren en $F$, pues es el vertice del tri\'angulo $\Delta GEF$
        \item Los puntos $K$, $L$ y $Q$, $T$ son colineales, de la misma manera $L$, $M$ y $V$, $S$ lo son, por tanto $\overline{QT}$ y $\overline{SV}$ concurren en $L$, pues es el vertice del tri\'angulo $\Delta LMK$
    \end{itemize}
    Utilizando el teorema de Desargues, tenemos que los tri\'angulos $\Delta OVS$ y $\Delta TQN$ est\'an en perspectiva desde un punto, por tanto como los vertices de ambos tri\'angulos estan distribuidos de forma correcta sobre las rectas $l_1$, $l_2$ y $l_3$ son concurrentes
\begin{figure}[h!]
    \centering
    \includegraphics[width=0.9\textwidth]{desargues.jpeg}
    \caption{Construcci\'on}
\end{figure}
\newpage
\textbf{Casos degenerados:}
\begin{itemize}
    \item Si los tres tri\'angulos tienen centro de perspectiva en el punto al infinito, se puede hacer una construcci\'on an\'aloga a la anterior
    \item Si los tres son tri\'angulos homoteticos, entonces los tres ejes de perspectiva coinciden, son la recta al infinito.
\end{itemize}
\begin{figure}[h!]
    \centering
    \includegraphics[width=0.9\textwidth]{Captura de pantalla (784).png}
    \caption{Casos degenerados}
\end{figure}
\end{proof}
\end{document}